% P15.26: Durrett 3.3.1--3.3.21, with the full Section 3.3 target retained.
\chapter{Characteristic Functions and Inversion}
\label{sec:p15-characteristic-functions-and-inversion}

A characteristic function represents a probability law by bounded
oscillating integrals. Multiplication describes independent addition,
inversion recovers the law, and continuity near zero controls escape
of mass. This chapter supplies the definitions, numbered theorems
and examples of the first three subsections of Durrett's
Section~3.3, through Theorem~3.3.21
\cite[pp.~125--136]{DurrettPTEAuthor2019}.
The section's later Polya and moment-problem presentations remain
separate required components, explicitly identified at the end.

\section{Complex-valued integration and differentiation}
\label{sec:p15-characteristic-complex-foundations}

\begin{hbdefinition}{def:p15-characteristic-complex-conventions}{Complex numbers and integrable complex functions}
Write \(z=a+ib\), \(a,b\in\mathbb R\), where \(i^2=-1\).
Its real and imaginary parts are \(\operatorname{Re}z=a\),
\(\operatorname{Im}z=b\), its conjugate is \(\overline z=a-ib\),
and \(|z|=(a^2+b^2)^{1/2}\).
Set \(e^{a+ib}=e^a(\cos b+i\sin b)\).
A complex-valued function is measurable when its two real
components are measurable. It is integrable against a measure
\(\nu\) when \(\int|f|\,d\nu<\infty\), and then
\[
 \int f\,d\nu=\int\operatorname{Re}f\,d\nu
                       +i\int\operatorname{Im}f\,d\nu.
\]
Complex expectations use this convention. Differentiation of a
complex function of a real parameter means differentiation of
both real components. Integrals \(\int_0^h\) are oriented
when \(h<0\).
\end{hbdefinition}

\begin{hblemma}{thm:p15-characteristic-complex-calculus}{proved=annex:p15.characteristic-complex-calculus}{The needed complex integral and exponential rules}
Complex integrability is equivalent to integrability of both real
components. Linearity and conjugation commute with integration, and
\[
 \left|\int f\,d\nu\right|\leq\int|f|\,d\nu.
\]
Dominated convergence holds with an integrable real bound on
the modulus. For sigma-finite product measures, Fubini holds
whenever the modulus is integrable.
For every \(z,w\in\mathbb C\),
\[
 e^z=\sum_{k=0}^\infty z^k/k!,\qquad
 e^{z+w}=e^ze^w,\qquad
 \frac{d}{dx}e^{zx}=ze^{zx}\quad(x\in\mathbb R).
\]
The series is absolutely convergent.
For real \(v\),
\[
 |e^{iv}|=1,\qquad |e^{iv}-1|\leq\min\{|v|,2\}.
\]
\end{hblemma}

\begin{hbtheorem}{thm:p15-characteristic-differentiation}{proved=annex:p15.characteristic-differentiation}{The complete differentiating-under-the-integral premise}
Let \((S,\mathcal A,\nu)\) be a measure space, \(y\in\mathbb R\),
\(\delta>0\), and \(I=(y-\delta,y+\delta)\).
Suppose \(f:I\times S\to\mathbb C\) satisfies:
for every \(t\in I\), \(f(t,\cdot)\) and
\(g(t,\cdot)=\partial_t f(t,\cdot)\) are measurable and
integrable; for every \(s\), \(f(\cdot,s)\) is continuously
differentiable on \(I\); the function
\(v(t)=\int g(t,s)\,\nu(ds)\) is continuous at \(y\); and
\[
 \int_S\int_I |g(t,s)|\,dt\,\nu(ds)<\infty.
\]
Then \(u(t)=\int f(t,s)\,\nu(ds)\) is differentiable at \(y\)
and \(u'(y)=v(y)\).
No sigma-finiteness of \(\nu\) is required: the proof reduces
the derivative to an explicitly constructed sigma-finite support.
This supplies the source Appendix A.5.1 argument
\cite[pp.~470--471]{DurrettPTEAuthor2019}.
\end{hbtheorem}

\section{Definition, algebra and the source examples}
\label{sec:p15-characteristic-definition-and-examples}

\begin{hbdefinition}{def:p15-characteristic-function}{The characteristic function of a real law}
For a finite real random variable \(X\) of Borel probability law
\(\mu\), define
\[
 \varphi_X(t)=\mathbb E e^{itX}
       =\int_{\mathbb R}e^{itx}\,\mu(dx),\qquad t\in\mathbb R.
\]
Thus \(\varphi_X(t)=\mathbb E\cos(tX)+i\mathbb E\sin(tX)\);
it depends only on the law. A function is called a
characteristic function only when it has this representation
for a Borel probability measure.
The symbol \(\phi(x)=(2\pi)^{-1/2}e^{-x^2/2}\) continues
to denote the standard normal density, not a characteristic
function.
\index{characteristic function}
\end{hbdefinition}

\begin{hbtheorem}{thm:p15-characteristic-algebra}{proved=annex:p15.characteristic-algebra}{Durrett 3.3.1, 3.3.2 and 3.3.9}
Every characteristic function satisfies
\[
 \varphi(0)=1,\quad \varphi(-t)=\overline{\varphi(t)},\quad
 |\varphi(t)|\leq1,\quad
 |\varphi(t+h)-\varphi(t)|\leq\mathbb E|e^{ihX}-1|.
\]
It is uniformly continuous on \(\mathbb R\).
For real \(a,b\), the law of \(aX+b\) has characteristic
function \(e^{itb}\varphi_X(at)\), including \(a=0\).
Independent \(X,Y\) have
\(\varphi_{X+Y}=\varphi_X\varphi_Y\).
For probability laws \(\mu_j\) and nonnegative weights
\(\lambda_j\) summing to one, the finite mixture
\(\sum_j\lambda_j\mu_j\) has characteristic function
\(\sum_j\lambda_j\varphi_j\).
The same assertion holds for a countable mixture.
Also \(\operatorname{Re}\varphi\) and \(|\varphi|^2\)
are characteristic functions: they arise from an equal mixture
of \(X,-X\), and from \(X-X'\) for an independent copy \(X'\).
\end{hbtheorem}

\begin{hbtheorem}{thm:p15-characteristic-elementary-examples}{proved=annex:p15.characteristic-elementary-examples}{The elementary source transforms}
The following probability laws have the indicated characteristic
functions, with all displayed removable values at \(t=0\) set to one.
The uniform endpoints \(a<b\) are finite real numbers.
\[
 \begin{array}{c|c}
 \text{law or density}&\text{characteristic function}\\ \hline
 \mathbb P(X=1)=\mathbb P(X=-1)=1/2 & \cos t\\
 \operatorname{Poisson}(\lambda),\ \lambda\geq0
                    & \exp\{\lambda(e^{it}-1)\}\\
 (b-a)^{-1}\mathbf1_{\{a<x<b\}},\ a<b
                    &(e^{itb}-e^{ita})/(it(b-a))\\
 (1-|x|)^+          &2(1-\cos t)/t^2\\
 e^{-x}\mathbf1_{\{x>0\}}&1/(1-it)\\
 \tfrac12e^{-|x|}   &1/(1+t^2)
 \end{array}
\]
Here \(r^+=\max\{r,0\}\), and the Poisson law at \(\lambda=0\)
is the point mass at zero.
For the uniform law on \((-c,c)\), \(c>0\), the transform is
\(\sin(ct)/(ct)\).
The triangular example is the sum of two independent uniforms
on \((-1/2,1/2)\). The bilateral exponential example is the
source's mixture of an exponential law and its reflection.
These are Examples~3.3.3--3.3.4, 3.3.6--3.3.8 and 3.3.10.
\end{hbtheorem}

\begin{hbtheorem}{thm:p15-characteristic-normal-transform}{proved=annex:p15.characteristic-normal-transform}{Durrett 3.3.5: both Gaussian calculations}
A standard normal variable has characteristic function
\(\exp(-t^2/2)\). A normal law of real mean \(m\) and finite variance
\(v\geq0\) has transform
\[
 \exp(imt-vt^2/2).
\]
At \(v=0\) the law is the point mass at \(m\).
Both the source's real differentiation/integration-by-parts
proof and its completed-square route are supplied. The latter
uses the rigorously established shifted integral
\[
 \frac1{\sqrt{2\pi}}\int_{\mathbb R}e^{-(x-is)^2/2}\,dx=1
                                      \quad(s\in\mathbb R);
\]
this is an integrable complex function, not a probability
density with a complex mean.
\end{hbtheorem}

\section{The oscillatory integral and inversion}
\label{sec:p15-characteristic-inversion}

\begin{hbdefinition}{def:p15-characteristic-inversion-kernels}{Removable values and the symmetric improper integral}
Set \(\operatorname{sinc}(v)=\sin(v)/v\) for \(v\ne0\),
and \(\operatorname{sinc}(0)=1\).
For \(A\geq0\), define
\(S(A)=\int_0^A\operatorname{sinc}(x)\,dx\).
For real \(\theta\) and \(T>0\), put
\[
 R(\theta,T)=\int_{-T}^T\frac{\sin(\theta t)}t\,dt,
\]
where the integrand at zero has value \(\theta\).
For finite real \(a<b\), define
\[
 K_{a,b}(t)=\int_a^b e^{-ity}\,dy
 =\frac{e^{-ita}-e^{-itb}}{it}\ (t\ne0),\qquad
 K_{a,b}(0)=b-a.
\]
Every inversion limit below means the limit of the finite
integrals over \([-T,T]\), as \(T\to\infty\). It does not
assume absolute integrability over the entire real line.
Write \(\operatorname{sgn}\theta=1,0,-1\) according as
\(\theta>0,\theta=0,\theta<0\).
\end{hbdefinition}

\begin{hblemma}{thm:p15-characteristic-dirichlet-integral}{proved=annex:p15.characteristic-dirichlet-integral}{The complete Dirichlet-integral premise}
For \(A>0\), the source strip-integration argument gives
\[
 S(A)=\frac\pi2-\cos A\int_0^\infty\frac{e^{-Ay}}{1+y^2}\,dy
                -\sin A\int_0^\infty\frac{y e^{-Ay}}{1+y^2}\,dy.
\]
In fact \(|S(A)-\pi/2|\leq A^{-1}\) and
\(|S(A)-\arctan A|\leq2/A\), so \(S(A)\to\pi/2\).
The function \(S\) is bounded in absolute value on \([0,\infty)\),
and
\[
 R(\theta,T)=2\operatorname{sgn}(\theta)S(T|\theta|)
           \longrightarrow\pi\operatorname{sgn}(\theta).
\]
The convergence of \(\int_0^\infty(\sin x)/x\,dx\) is not
absolute. All the finite-strip interchanges are nevertheless
absolutely integrable.
\end{hblemma}

\begin{hbtheorem}{thm:p15-characteristic-inversion-uniqueness}{proved=annex:p15.characteristic-inversion-uniqueness}{Durrett 3.3.11: inversion, atoms and uniqueness}
If \(\varphi(t)=\int e^{itx}\,\mu(dx)\) for a Borel probability
measure \(\mu\), then for every finite \(a<b\),
\[
 \lim_{T\to\infty}\frac1{2\pi}\int_{-T}^T
                       K_{a,b}(t)\varphi(t)\,dt
 =\mu((a,b))+\tfrac12\mu(\{a\})+\tfrac12\mu(\{b\}).
\]
The limit exists as stated. Two Borel probability laws with
the same characteristic function are equal.
For \(\mu=\delta_0\), \(a=-1,b=1\), the integrand is
\(2\sin t/t\), giving an explicit case without absolute
convergence.
For the same law with \(a=0,b=1\), the limit is \(1/2\)
although \(\mu((0,1))=0\), exhibiting the endpoint correction.
\index{inversion formula!characteristic functions}
\end{hbtheorem}

\begin{hbcorollary}{thm:p15-characteristic-inversion-corollaries}{proved=annex:p15.characteristic-inversion-corollaries}{Durrett 3.3.12--3.3.13}
A real law is symmetric about zero if and only if its
characteristic function is real-valued.
If independent centered normal variables have variances
\(v_1,v_2\geq0\), their sum is centered normal with variance
\(v_1+v_2\).
The proof here uses characteristic functions and uniqueness,
retaining this source route separately from the earlier
density-convolution calculation.
\end{hbcorollary}

\begin{hbtheorem}{thm:p15-characteristic-integrable-inversion}{proved=annex:p15.characteristic-integrable-inversion}{Durrett 3.3.14: an integrable transform gives a density}
Let \(\varphi\) be the characteristic function of the Borel
probability law \(\mu\).
If \(\int_{\mathbb R}|\varphi(t)|\,dt<\infty\), its law
has the bounded, uniformly continuous density
\[
 f(y)=\frac1{2\pi}\int_{\mathbb R}e^{-ity}\varphi(t)\,dt,
 \qquad
 0\leq f(y)\leq\frac1{2\pi}\int_{\mathbb R}|\varphi(t)|\,dt.
\]
In particular the law has no atoms and \(\int f(y)\,dy=1\).
The proof establishes reality, nonnegativity and normalization
of the inverse integral, not just its formal expression.
\end{hbtheorem}

\begin{hbtheorem}{thm:p15-characteristic-inverse-examples}{proved=annex:p15.characteristic-inverse-examples}{Durrett 3.3.15--3.3.16 and the Cauchy-sum consequence}
The density
\[
 f_P(x)=\frac{1-\cos x}{\pi x^2}\quad(x\ne0),\qquad
 f_P(0)=\frac1{2\pi},
\]
has characteristic function \((1-|t|)^+\).
This is the source's Polya distribution.
The standard Cauchy density \(1/(\pi(1+x^2))\) has
characteristic function \(e^{-|t|}\).
Both identities follow by applying integrable inversion to
the preceding triangular and bilateral exponential examples.
If \(X_1,\ldots,X_n\) are independent standard Cauchy variables,
then \((X_1+\cdots+X_n)/n\) has the same law as \(X_1\).
This supplies the separate characteristic-function proof
indicated in source Exercise~3.3.6.
\end{hbtheorem}

\begin{hbtheorem}{thm:p15-characteristic-atom-inversion}{proved=annex:p15.characteristic-atom-inversion}{Atom recovery and the lattice formulas}
For a real probability law with characteristic function \(\varphi\),
\[
 \mu(\{a\})=\lim_{T\to\infty}\frac1{2T}
                   \int_{-T}^T e^{-ita}\varphi(t)\,dt.
\]
If \(\mathbb P(X\in b+h\mathbb Z)=1\), where \(h>0,b\in\mathbb R\),
then, for \(x\in b+h\mathbb Z\),
\[
 \mathbb P(X=x)=\frac h{2\pi}
          \int_{-\pi/h}^{\pi/h}e^{-itx}\varphi(t)\,dt,\qquad
 \varphi(t+2\pi/h)=e^{2\pi ib/h}\varphi(t).
\]
For \(b=0\) this is ordinary periodicity.
Also
\[
 \lim_{T\to\infty}\frac1{2T}\int_{-T}^T|\varphi(t)|^2\,dt
          =\sum_{x:\mu(\{x\})>0}\mu(\{x\})^2.
\]
In particular \(\varphi(t)\to0\) as \(t\to+\infty\) implies
that the law has no atoms. These are the inversion consequences
in Exercises~3.3.2--3.3.3.
\end{hbtheorem}

\section{Weak convergence and characteristic functions}
\label{sec:p15-characteristic-continuity}

\begin{hbtheorem}{thm:p15-characteristic-continuity-theorem}{proved=annex:p15.characteristic-continuity-theorem}{Durrett 3.3.17: the complete continuity theorem}
Let \(\mu_n\) be Borel probability measures with characteristic
functions \(\varphi_n\).
If \(\mu_n\Rightarrow\mu_\infty\), then
\(\varphi_n(t)\to\varphi_\infty(t)\) for every real \(t\).
Conversely, if \(\varphi_n(t)\to\varphi(t)\) for every real \(t\)
and the limit is continuous at zero, then the laws are tight
and converge weakly to a unique probability law with
characteristic function \(\varphi\).

The source's tightness estimate is, for \(u>0\),
\[
 \mu_n(\{|x|\geq2/u\})
 \leq\frac1u\int_{-u}^u(1-\operatorname{Re}\varphi_n(t))\,dt
 =2\int_{\mathbb R}(1-\operatorname{sinc}(ux))\,\mu_n(dx).
\]
The middle expression also equals the integral with
\(1-\varphi_n(t)\); its imaginary part cancels, so it is
a real nonnegative quantity.
The source's bounded-test subsequence argument is included
in full.
\end{hbtheorem}

\section{Moments, derivatives and the Taylor remainder}
\label{sec:p15-characteristic-moments}

In this section \(X\) is a finite real random variable, \(\mu\)
is its Borel probability law, and
\(\varphi(t)=\int e^{itx}\,\mu(dx)\).

\begin{hbtheorem}{thm:p15-characteristic-moment-derivatives}{proved=annex:p15.characteristic-moment-derivatives}{Durrett 3.3.18}
If \(n\geq1\) is an integer and \(\mathbb E|X|^n<\infty\),
then \(\varphi\) has continuous derivatives through order \(n\),
and
\[
 \varphi^{(j)}(t)=\mathbb E[(iX)^j e^{itX}]
                 \quad(0\leq j\leq n,\ t\in\mathbb R).
\]
In particular \(\varphi^{(j)}(0)=i^j\mathbb EX^j\).
Every lower absolute moment is finite, and the derivative
interchanges are justified individually.
\end{hbtheorem}

\begin{hblemma}{thm:p15-characteristic-taylor-remainder}{proved=annex:p15.characteristic-taylor-remainder}{Durrett 3.3.19: both remainder bounds}
For every integer \(n\geq0\) and every real \(x\),
\[
 \left|e^{ix}-\sum_{m=0}^n\frac{(ix)^m}{m!}\right|
 \leq\min\left\{\frac{|x|^{n+1}}{(n+1)!},
                         \frac{2|x|^n}{n!}\right\}.
\]
Zeroth powers denote the constant one, including at zero.
The exact oriented-integral remainder is
\[
 e^{ix}-\sum_{m=0}^n\frac{(ix)^m}{m!}
       =\frac{i^{n+1}}{n!}\int_0^x(x-s)^n e^{is}\,ds.
\]
The case \(n=0\) is covered separately where the second
integration-by-parts argument would involve \((n-1)!\).
\end{hblemma}

\begin{hbtheorem}{thm:p15-characteristic-moment-expansion}{proved=annex:p15.characteristic-moment-expansion}{The full moment expansion and Durrett 3.3.20}
For an integer \(n\geq0\) with \(\mathbb E|X|^n<\infty\),
\[
 \varphi(t)=\sum_{m=0}^n\frac{(it)^m\mathbb EX^m}{m!}
                                      +o(|t|^n).
\]
Here \(o(|t|^n)\) is a complex remainder whose modulus divided
by \(|t|^n\) tends to zero; \(X^0=1\).
At every real \(t\), its absolute error is at most
\[
 \mathbb E\min\left\{\frac{|tX|^{n+1}}{(n+1)!},
                              \frac{2|tX|^n}{n!}\right\}
 \leq\mathbb E\min\{|tX|^{n+1},2|tX|^n\}.
\]
Thus if \(\mathbb EX^2<\infty\),
\[
 \varphi(t)=1+it\mathbb EX-\tfrac12t^2\mathbb EX^2+o(t^2),
\]
with no third-moment hypothesis.
\end{hbtheorem}

\begin{hbtheorem}{thm:p15-characteristic-second-moment}{proved=annex:p15.characteristic-second-moment}{Durrett 3.3.21 and the zero-law consequence}
The symmetric quotient
\[
 Q(h)=\frac{\varphi(h)-2\varphi(0)+\varphi(-h)}{h^2}
                    \quad(h\ne0)
\]
is real and nonpositive. One has
\[
 \mathbb EX^2<\infty
 \quad\Longleftrightarrow\quad
 \limsup_{h\downarrow0}Q(h)>-\infty.
\]
When this holds, \(Q(h)\to-\mathbb EX^2\) as \(h\to0\).
If the finite complex limit
\(c=\lim_{t\downarrow0}(\varphi(t)-1)/t^2\) exists, then
\(\mathbb EX=0\), \(c\) is real, and \(\mathbb EX^2=-2c\).
In particular \(\varphi(t)=1+o(t^2)\) implies \(X=0\)
almost surely and \(\varphi\equiv1\).
\end{hbtheorem}

\begin{hbtheorem}{thm:p15-characteristic-boundaries}{proved=annex:p15.characteristic-boundaries}{Explicit limits of the hypotheses}
The centered normal laws of variance \(n\) have transforms
\(e^{-nt^2/2}\), tending pointwise to \(\mathbf1_{\{0\}}(t)\).
This limit is discontinuous at zero; the laws are not tight
and have defective CDF limit \(1/2\).
Having a density does not imply an integrable characteristic
function: both a nondegenerate uniform law and a rate-one
exponential law give counterexamples.
The product rule requires independence: for one fair sign
\(X=Y\), the sum has transform \(\cos(2t)\), whereas the
product of the two marginal transforms is \(\cos^2t\).
The function \(e^{-|t|^\alpha}\) is not a characteristic
function when \(\alpha>2\).
\end{hbtheorem}

\begin{proofboundary}{def:p15-characteristic-next-boundary}{The full Section 3.3 target and the next exact chain}
This chapter supplies the source theorems, lemmas, corollaries
and examples numbered 3.3.1--3.3.21, and the stated additional
inversion consequences. The complete remaining Section~3.3 chain is
supplied in \cref{sec:p15-characteristic-function-criteria-and-moments}:
Polya's criterion, both stable-function routes, the periodic Fourier
example, all three moment counterexamples, the analytic and Carleman
criteria, and the complete method of moments. Together the two
chapters supply every Exercise~3.3.1--3.3.25.
The central-limit constructions and all Section~3.4 exercises are
supplied in \cref{sec:p15-central-limit-theorems}. The complete
normal-domain theorem and its normalization construction are
\cref{thm:p15-nda-characterization,thm:p15-nda-normalizer-construction};
the original cutoff route is \cref{thm:p15-nda-gk-cutoff-route}.
The later full stable-law presentations retain their distinct source
constructions.
\end{proofboundary}
